\documentclass[12pt]{article}
\usepackage[utf8]{inputenc}
\usepackage[T1]{fontenc}
\usepackage{hyperref}
\title{Fa d'Ambu Language Resources for K-12 Curriculum Development}
\author{Compiled by Claude Sonnet 5.5 (claude-sonnet-5-5) for Andr\'as Kornai}
\date{}
\begin{document}
\maketitle

\section{Introduction}
Fa d'Ambu (Fa d'Amb\^o, Annobonese Creole; ISO 639-3 \texttt{fab}) is a Portuguese-based creole of Annob\'on island, Equatorial Guinea, related to the creoles of S\~ao Tom\'e and Pr\'incipe. Estimates of speakers range from about 5,000 to 6,600. It is recognized as a minority language but has no education policy or standardization; Spanish is the official language.

\section{Primary Educational Resources}
No school materials were found in the search results.

\section{Linguistic and Reference Materials}
\begin{itemize}
\item Annobonese Creole (Wikipedia): \url{https://en.wikipedia.org/wiki/Annobonese_Creole}
\item Fa d'Amb\^o: from past to present (research article): \url{https://oasisbr.ibict.br/vufind/Record/RCAP_e3b38b5d95bacde146ac84a121bfd956}
\item Language Documentation and Description journal article (field-based snapshot): \url{https://www.lddjournal.org/article/id/1262}
\item Work Papers of the Summer Institute of Linguistics, University of North Dakota Session: \url{https://commons.und.edu/sil-work-papers/vol54/iss1/2}
\item De Gruyter article (International Journal of the Sociology of Language): \url{https://www.degruyterbrill.com/document/doi/10.1515/ijsl-2016-0009/pdf}
\item CLUL bibliographic reference: \url{https://clul.ulisboa.pt/bibcite/reference/8689}
\end{itemize}
The exact subject of the last four items was not verified beyond the search listing.

\section{Religious and Community Resources}
None found.

\section{Digital and Audio Resources}
\begin{itemize}
\item Rosetta Panglossia entry for fab: \url{https://rosettapanglossia.longnow.org/index.php?title=Fab}
\end{itemize}

\section{Curriculum Development Context}
The language is still acquired by children, but there is no standard orthography or school use. Research literature is the main resource base; curriculum work would need new materials and community input.
\end{document}